% ─────────────────────────────────────────────────────────────────────────────
\chapter{Conclusions, Limitations, and Future Work}
\label{chap:conclusions}
% ─────────────────────────────────────────────────────────────────────────────

% ─────────────────────────────────────────────────────────────────────────────
\section{Summary of the Research}
\label{sec:conclusions-summary}
% ─────────────────────────────────────────────────────────────────────────────

This work originates from the need to systematically evaluate the operational
reliability of \gls{ai}-based systems used for forensic image triage. In many
application scenarios, these systems are still evaluated primarily on clean,
balanced, and semantically coherent data. In \gls{df} workflows, however, images
may exhibit much more heterogeneous characteristics: they may be ambiguous,
out-of-distribution, recompressed, resized, altered, or deliberately
manipulated. In this context, nominal performance alone is not sufficient to
establish whether a system can be reliably used to support investigative
review.

This thesis addressed this problem by adopting a perspective centered on
operational robustness. Adversarial attacks and anti-forensic transformations
were not treated as an autonomous objective of \gls{advml}, but as experimental
stressors useful for observing the behavior of proxy models and selected
software systems under non-ideal conditions. The main objective was therefore
to verify whether \gls{ai}-based automation exhibits coherent behavior when
exposed to clean inputs, \gls{ood} samples, anti-forensic transformations, and
adversarial perturbations. This analysis also included the observation and
normalization of the outputs produced by selected black-box software systems,
in order to compare the behavior of proxy models with that of systems used in
operationally relevant forensic-context scenarios.

The experimental process included the construction of a frozen dataset composed
of 1,500 images, the evaluation of the EfficientNet-B0, ResNet18, and
\gls{clip} proxy models under clean, \gls{ood}, anti-forensic, and adversarial
conditions, and the production of an 11,500-file forensic evaluation bundle
for the black-box evaluation of \gls{magnetaxiomtool}/\gls{magnetai},
\gls{excirefoto2025}, \gls{cellebriteinseyets}, and
\gls{griffeye}/\gls{t3kcore}.

The adopted experimental protocol is referred to as \gls{fairlab}. It does not
represent a new software tool or a new general architecture, but a reproducible
and auditable procedure for evaluating the operational robustness of
\gls{ai}-based systems in \gls{df} scenarios. Its structure makes it possible
to connect datasets, proxy models, perturbations, transformations, the blind
bundle, manifests, hashes, normalized black-box outputs, quantitative metrics,
and qualitative \gls{xai} analysis within a single traceable pipeline. In this
way, the coherence among dataset construction, stress testing, and black-box
evaluation does not depend on implicit steps, but on verifiable and
reconstructible artifacts.

The overall finding of this thesis is that the validation of \gls{ai}-based
systems in forensic contexts cannot be reduced to the verification of nominal
performance on clean data. Rather, it must be treated as an experimental
procedure specific to the task, system version, operational configuration, and
input conditions. From this perspective, robustness is not an abstract property
of the model, but an observable property of the system within the investigative
workflow in which it is used.

% ─────────────────────────────────────────────────────────────────────────────
\subsection{Excire Foto 2025 Results and Sensitivity Analysis}
\label{subsec:results-excire}
% ─────────────────────────────────────────────────────────────────────────────

The evaluation of \gls{excirefoto2025} requires specific interpretive caution.
The software is not treated as a product developed exclusively for forensic
purposes, nor as a native classifier for the binary
\texttt{weapon}/\texttt{non_weapon} distinction. In the experimental protocol,
it is evaluated as a black-box semantic retrieval system equipped with
\gls{ai}-based functionality. The operational prediction is therefore derived
from the presence or absence of a sample among the results retrieved by the
configured semantic search. This recoding enables a quantitative evaluation
against the experimental ground truth, but it must not be interpreted as access
to the internal classification logic of the system.

To evaluate the sensitivity of the operational behavior to the retrieval
threshold, three configurations of the \texttt{Distance limit} parameter were
considered: \texttt{D20}, \texttt{D50}, and \texttt{D80}. The \texttt{D20}
configuration represents a restrictive setting, oriented toward reducing false
positives and semantically weak retrievals; \texttt{D50} constitutes the
intermediate configuration of the sensitivity analysis; \texttt{D80} instead
represents a permissive setting, oriented toward maximizing recall but at the
cost of a substantial increase in false positives and in the weapon flag rate on
\gls{ood} samples.

\begin{table}[H]
\centering
\scriptsize
\renewcommand{\arraystretch}{1.15}
\setlength{\tabcolsep}{4pt}
\caption{Sensitivity analysis of \gls{excirefoto2025} with respect to the \texttt{Distance limit} parameter.}
\label{tab:results-excire-sensitivity}
\begin{tabularx}{0.95\textwidth}{@{}l r r r r r X@{}}
\toprule
\textbf{Setting} &
\textbf{Clean Acc.} &
\textbf{Recall} &
\textbf{\gls{fnr}} &
\textbf{\gls{fpr}} &
\textbf{OOD W-rate} &
\textbf{Operational profile} \
\midrule

\texttt{D20} &
0.927 &
0.880 &
0.120 &
0.026 &
0.238 &
Restrictive: low operational noise and lower \gls{ood} absorption, but higher
risk of false negatives. \

\texttt{D50} &
0.944 &
0.958 &
0.042 &
0.070 &
0.340 &
Intermediate: balanced compromise in the experimental context considered. \

\texttt{D80} &
0.915 &
0.986 &
0.014 &
0.156 &
0.522 &
Permissive: very high recall, but marked increase in false positives and
\gls{ood} flags. \

\bottomrule
\end{tabularx}
\end{table}

The comparison among the three configurations shows that the
\texttt{Distance limit} parameter does not act merely as a technical variable,
but as an operational threshold that modifies the risk profile of the system.
\texttt{D20} is more conservative and reduces the review workload arising from
false positives, but may fail to retrieve relevant \texttt{weapon} content.
\texttt{D50} represents, in the experimental context considered, the most
balanced configuration for triage, since it maintains high recall on the
\texttt{weapon} class without producing the level of noise observed under
\texttt{D80}. The latter strongly reduces false negatives, but substantially
increases false positives and maps more than half of the \gls{ood} inputs to
the \texttt{weapon} class.

\begin{table}[H]
\centering
\caption{Aggregate performance of \gls{excirefoto2025} under perturbed conditions.}
\label{tab:results-excire-perturbed-summary}
\small
\begin{tabular}{l l r r r r}
\toprule
\textbf{Setting} &
\textbf{Condition} &
\textbf{Accuracy} &
\textbf{Recall weapon} &
\textbf{\gls{fnr}} &
\textbf{\gls{fpr}} \
\midrule

\texttt{D20} & Anti-forensic & 0.927 & 0.875 & 0.125 & 0.022 \
\texttt{D50} & Anti-forensic & 0.941 & 0.957 & 0.043 & 0.074 \
\texttt{D80} & Anti-forensic & 0.908 & 0.984 & 0.016 & 0.169 \

\midrule

\texttt{D20} & Adversarial & 0.892 & 0.834 & 0.166 & 0.051 \
\texttt{D50} & Adversarial & 0.904 & 0.938 & 0.062 & 0.130 \
\texttt{D80} & Adversarial & 0.861 & 0.977 & 0.023 & 0.256 \

\bottomrule
\end{tabular}
\end{table}

The aggregate results under perturbed conditions show that, for
\gls{excirefoto2025}, anti-forensic transformations produce a more limited
impact than adversarial perturbations. However, in this case as well, the
operational profile changes substantially as the threshold varies:
\texttt{D20} maintains the lowest \gls{fpr} but exhibits the highest \gls{fnr},
\texttt{D80} almost completely reduces false negatives but increases
operational noise, whereas \texttt{D50} maintains an intermediate compromise
between recall and selectivity.

Table~\ref{tab:results-excire-d50-attack-metrics} reports the detailed results
for the \texttt{D50} configuration, which is used as the main reference for the
operational comparison. This choice does not imply that \texttt{D50} is optimal
in an absolute sense, but makes it possible to analyze an intermediate
configuration between the restrictive and permissive profiles.

\begin{table}[H]
\centering
\caption{Performance of \gls{excirefoto2025} \texttt{D50} under adversarial and anti-forensic conditions.}
\label{tab:results-excire-d50-attack-metrics}
\scriptsize
\renewcommand{\arraystretch}{1.12}
\setlength{\tabcolsep}{4pt}
\begin{tabular}{@{}l l r r r r r r@{}}
\toprule
\textbf{Family} &
\textbf{Condition} &
\textbf{Accuracy} &
\textbf{Recall} &
\textbf{\gls{fnr}} &
\textbf{\gls{fpr}} &
\textbf{FN} &
\textbf{FP} \
\midrule

Adversarial &
\gls{colorshift} &
0.949 &
0.966 &
0.034 &
0.068 &
17 &
34 \

Adversarial &
\gls{fgsm} &
0.926 &
0.890 &
0.110 &
0.038 &
55 &
19 \

Adversarial &
\gls{onepixel} &
0.946 &
0.964 &
0.036 &
0.072 &
18 &
36 \

Adversarial &
\gls{sigmazero} &
0.742 &
0.926 &
0.074 &
0.442 &
37 &
221 \

Adversarial &
\gls{superdeepfool} &
0.956 &
0.944 &
0.056 &
0.032 &
28 &
16 \

\midrule

Anti-forensic &
\gls{contraststretching} &
0.941 &
0.952 &
0.048 &
0.070 &
24 &
35 \

Anti-forensic &
\gls{gaussianblur} &
0.941 &
0.960 &
0.040 &
0.078 &
20 &
39 \

Anti-forensic &
\gls{histogrammodification} &
0.932 &
0.952 &
0.048 &
0.088 &
24 &
44 \

Anti-forensic &
\gls{jpegrecompression} &
0.946 &
0.958 &
0.042 &
0.066 &
21 &
33 \

Anti-forensic &
\gls{resampleresize} &
0.947 &
0.964 &
0.036 &
0.070 &
18 &
35 \

\bottomrule
\end{tabular}
\end{table}

A particularly relevant case is that of \gls{sigmazero}. Unlike a simple
evasion dynamic, in which the main effect would be the increase in false
negatives on the \texttt{weapon} class, \gls{excirefoto2025} shows, under the
\texttt{D50} and \texttt{D80} configurations, a semantic over-triggering
behavior. Under these conditions, the system continues to retrieve many
\texttt{weapon} images, but substantially increases the proportion of
\texttt{non_weapon} images erroneously retrieved as relevant.

\begin{table}[H]
\centering
\caption{Behavior of \gls{excirefoto2025} on the \gls{sigmazero} perturbation.}
\label{tab:results-excire-sigma-zero}
\small
\begin{tabular}{l r r r r r r}
\toprule
\textbf{Setting} &
\textbf{Accuracy} &
\textbf{Recall} &
\textbf{\gls{fnr}} &
\textbf{\gls{fpr}} &
\textbf{FN} &
\textbf{FP} \
\midrule

\texttt{D20} & 0.815 & 0.820 & 0.180 & 0.190 & 90 & 95 \
\texttt{D50} & 0.742 & 0.926 & 0.074 & 0.442 & 37 & 221 \
\texttt{D80} & 0.657 & 0.976 & 0.024 & 0.662 & 12 & 331 \

\bottomrule
\end{tabular}
\end{table}

The joint interpretation of the three configurations shows that
\gls{sigmazero} does not give rise only to a missed-detection problem. On the
contrary, as the retrieval threshold increases, recall on the \texttt{weapon}
class also increases, but the \gls{fpr} grows markedly as well: from 0.190 under
\texttt{D20}, to 0.442 under \texttt{D50}, and up to 0.662 under \texttt{D80}.
This behavior is operationally significant because, while a permissive
configuration may appear preferable due to its lower risk of false negatives,
it can generate a high volume of false positives, saturating human review and
reducing the efficiency of the automatic filter.

Overall, \gls{excirefoto2025} highlights the critical role of the retrieval
threshold in black-box evaluation. A restrictive configuration reduces noise but
increases the risk of missed detection; a permissive configuration increases
recall but may strongly amplify operational noise, especially in the presence
of perturbed or semantically ambiguous inputs. For this reason, in the
subsequent comparative discussion, the \texttt{D50} configuration is used as
the main reference, whereas \texttt{D20} and \texttt{D80} are interpreted as
operational extremes of the sensitivity analysis.

% ─────────────────────────────────────────────────────────────────────────────
\section{Main Experimental Findings}
\label{sec:conclusions-findings}
% ─────────────────────────────────────────────────────────────────────────────

The experimental results confirm that operational robustness cannot be evaluated
solely on the basis of the clean baseline. Nominal performance on clean data is
a necessary starting point, but it is not sufficient to estimate the reliability
of an \gls{ai}-based system in a forensic workflow. In the analyzed context,
proxy models and selected software systems with high aggregate performance may
still exhibit critical behaviors when exposed to out-of-distribution inputs,
anti-forensic transformations, or adversarial perturbations.

This evidence is particularly relevant because forensic triage does not operate
on ideal data. Images acquired during an investigation may be recompressed,
resized, partially degraded, semantically ambiguous, or deliberately
manipulated. For this reason, the experimental evaluation considered multiple
input conditions, observing not only the ability of the systems to correctly
classify clean images, but also their stability with respect to operationally
plausible or adversarially manipulated conditions.
Table~\ref{tab:conclusions-main-findings} summarizes the main findings, linking
each experimental area to its corresponding operational implication.

\begin{table}[H]
\centering
\caption{Interpretive summary of the main experimental findings.}
\label{tab:conclusions-main-findings}
\footnotesize
\renewcommand{\arraystretch}{1.2}
\begin{tabularx}{\textwidth}{@{}p{0.18\textwidth} X X@{}}
\toprule
\textbf{Area} & \textbf{Main finding} & \textbf{Operational relevance} \
\midrule
Clean baseline
& EfficientNet-B0 shows the highest clean performance among the proxy models
(0.978); among the selected software systems, \gls{griffeye}/\gls{t3kcore}
achieves the highest clean accuracy under the adopted conservative mapping,
whereas \gls{cellebriteinseyets} exhibits a balanced profile with high
\texttt{weapon} recall.
& Nominal performance is necessary to describe the baseline behavior of the
system, but it is not sufficient to evaluate its reliability under operational
conditions. \

\gls{ood}
& \gls{clip} exhibits the highest rate of \texttt{weapon} predictions on
\gls{ood} samples (0.8356); EfficientNet-B0 and ResNet18 instead exhibit high
confidence levels even in the presence of inputs outside the nominal binary
task.
& Out-of-distribution inputs may be absorbed into the \texttt{weapon} class,
generating operational noise, misleading prioritization, or additional
workload for the analyst. \

Anti-forensic
& Anti-forensic transformations generally produce limited drops, but introduce
operationally relevant errors, especially false negatives on
\texttt{weapon}.
& Operationally plausible manipulations such as blurring, recompression,
resizing, or tonal alterations may affect automated triage even when the
aggregate metric degradation appears limited. \

Adversarial
& EfficientNet-B0 shows the most marked degradation under \gls{sigmazero}
(accuracy 0.015); in selected black-box software systems, the effects of
perturbations are heterogeneous and depend on the system, configuration, and
operational mapping.
& High clean accuracy does not guarantee robustness against deliberately
manipulated inputs, especially when the system is exposed to perturbations
designed to alter its decision. \

Selected software systems
& The selected software systems show differentiated operational profiles:
\gls{griffeye}/\gls{t3kcore} is more selective under the adopted conservative
mapping, \gls{cellebriteinseyets} maintains a balanced profile, and
\gls{magnetai} and \gls{excirefoto2025} exhibit specific failure modes.
& Selected software systems must be interpreted through observable outputs,
operational thresholds, mapping rules, and black-box normalization, avoiding
uncontrolled direct comparisons with proxy models. \

\gls{xai}
& \gls{ig} maps make representative cases directly observable, but do not
constitute stand-alone evidence of robustness.
& \gls{xai} supports critical review and qualitative diagnosis of failure modes,
without replacing quantitative metrics or expert validation. \
\bottomrule
\end{tabularx}
\end{table}

The joint interpretation of these results reveals a significant dissociation
between nominal accuracy and forensic reliability. Under clean conditions, some
systems achieve high performance and appear suitable for the considered task;
however, the same behavior is not necessarily preserved when the input
distribution, the visual quality of the image, or the nature of the manipulation
changes. This aspect is central in an investigative context, because errors do
not always have the same operational weight: a false negative on
\texttt{weapon} may reduce the probability that relevant content is submitted to
human review, whereas a false positive or an improper flag on \gls{ood} samples
may increase operational noise and alter the prioritization of the material to
be analyzed.

The results on \gls{ood} samples reinforce this interpretation. The fact that
some systems assign \texttt{weapon} labels or high confidence levels even to
images that are outside the nominal task indicates that automation may produce
apparently certain decisions in the absence of investigatively relevant content.
This behavior is not simply a classification error in the traditional sense,
but an operational risk related to uncertainty management and to the system's
ability to recognize the limits of its application domain.

Anti-forensic transformations show a less extreme effect than the most
aggressive adversarial attacks, but remain forensically relevant. \gls{jpeg}
recompression, blurring, resizing, and tonal modifications are operationally
plausible image-level transformations and may be introduced either by ordinary
sharing processes or by conduct deliberately aimed at reducing the automatic
readability of the content. Even when the aggregate accuracy drop remains
limited, the emergence of false negatives on \texttt{weapon} images shows that
robustness must be interpreted with respect to its effects on the investigative
workflow, not only with respect to the average variation in the metrics.

Adversarial attacks instead highlight the clearest limitation of nominal
performance. The case of EfficientNet-B0 under \gls{sigmazero} shows that a
model that is highly accurate under clean conditions may collapse in the
presence of perturbations constructed to alter its decision. For selected
black-box software systems, the interpretation is necessarily more cautious,
since their architectures, internal thresholds, preprocessing strategies, and
retrieval or categorization criteria are not known. For this reason, the
comparison must be based on observable outputs and on their operational
recoding, avoiding the attribution of unverifiable internal properties to the
systems.

Overall, the main experimental finding of this thesis is that validation on
clean data is not equivalent to validation under operational conditions. A
system may be accurate at baseline and, at the same time, produce false
negatives on \texttt{weapon} images, improper flags on \gls{ood} samples, or
retrieval-threshold-dependent errors. Operational robustness must therefore be
evaluated as a contextual property of the system within the forensic workflow,
jointly considering clean performance, \gls{ood} behavior, sensitivity to
manipulations, black-box outputs, and the need for human review.

% ─────────────────────────────────────────────────────────────────────────────
\section{Operational and Forensic Implications}
\label{sec:conclusions-operational-implications}
% ─────────────────────────────────────────────────────────────────────────────

The experimental evidence has direct implications for the responsible use of
\gls{ai}-based systems in forensic workflows. The central issue concerns the
role that such systems can play in investigative triage. Automated
classification can support evidence prioritization and reduce the initial
workload of the analyst, but its output must be interpreted as an investigative
indication to be verified, contextualized, and documented, not as an autonomous
evidentiary determination. Its usefulness depends on the possibility of
reconstructing the context in which the result was generated, the system
configuration, the reference dataset, the adopted operational thresholds, and
the subsequent human review. This requirement is particularly relevant for
selected black-box software systems, in which the observability of the
decision-making process is limited to exported outputs and their experimental
normalization.

From this perspective, false negatives take on specific operational relevance.
In a forensic triage workflow, a false negative on the \texttt{weapon} class may
be more critical than a false positive, because it may reduce the probability
that investigatively relevant content is submitted to human review. This risk
does not concern only proxy models, but also selected software systems, and may
increase in the presence of perturbations or transformations capable of
altering the automated readability of image content. The most severe case
observed in the protocol is that of EfficientNet-B0 under \gls{sigmazero}, with
489 \texttt{weapon}(\rightarrow)\texttt{non_weapon} transitions. This result
clearly shows that a system that is accurate under clean conditions may become
operationally fragile when exposed to deliberately manipulated inputs.

An additional element of caution concerns \gls{ood} inputs. In this work, these
samples were not treated as a conventional third supervised class, but as an
autonomous operational risk. Out-of-distribution images may be improperly
absorbed into the \texttt{weapon} class, generate non-relevant flags, increase
the analyst's review workload, or produce misleading prioritization. In forensic
contexts, this behavior is particularly concerning, because an apparently
confident output on content outside the task domain may induce a false
perception of system reliability.

The operational configuration of semantic retrieval systems represents another
factor that directly affects the interpretation of results. In the case of
\gls{excirefoto2025}, the \texttt{Distance limit} parameter modifies the
trade-off between false negatives and false positives: the \texttt{D20}
configuration is more restrictive, \texttt{D80} is more permissive, and
\texttt{D50} constitutes an intermediate compromise in the experimental context.
It follows that system configuration cannot be considered a secondary technical
detail, but must be documented as an integral part of the forensic workflow.
Given the same analyzed content, different thresholds may produce different
results and affect the subsequent review process.

These considerations reinforce the importance of traceability and auditability
throughout the entire process. The \gls{fairlab} pipeline keeps datasets,
manifests, hashes, the blind bundle, predictions, metrics, and system outputs
separate, making it possible to reconstruct the experimental path and to
distinguish among classification errors, missed detections, non-interpretable
outputs, and limitations of the operational mapping. This separation not only
has technical value, but also constitutes an essential methodological condition
for making the forensic use of \gls{ai} verifiable and defensible. In the
absence of traceable artifacts, the result produced by an automated system risks
remaining an opaque indication that is difficult to verify and poorly
contextualized.

Overall, the results confirm the need to maintain a
\textit{human-in-the-loop} paradigm when using \gls{ai}-based systems in
forensic contexts. Automated classification should be understood as support for
investigative triage, to be integrated with human review, documentation of the
operations, and contextualization of the result. This approach is consistent
with the principles of human oversight, reliability, and accountability set out
in the European guidelines for trustworthy \gls{ai}~\cite{EUAI2020}. This
requirement is particularly important for content of high investigative
relevance, out-of-distribution inputs, manipulated images, and outputs produced
by selected black-box software systems; in all these cases, operational
reliability cannot be presumed, but must be evaluated, documented, and subjected
to critical scrutiny.


% ─────────────────────────────────────────────────────────────────────────────
\section{Limitations}
\label{sec:conclusions-limitations}
% ─────────────────────────────────────────────────────────────────────────────

This work presents several limitations that must be explicitly stated so that
the conclusions can be interpreted correctly. These limitations do not reduce
the validity of the experimental protocol, but delimit its scope of application
and clarify the conditions under which the results can be considered
representative.

The main interpretive constraint concerns the dataset and the experimental
task. The constructed corpus is aimed at a specific investigative problem,
namely binary classification between \texttt{weapon} and \texttt{non_weapon},
with a separate evaluation of \gls{ood} samples. The overall size, equal to
1,500 images, enables a controlled and traceable evaluation, but does not cover
the full variability of visual content that may arise in real investigations.
For this reason, the results must not be automatically generalized to
multi-class tasks, different visual domains, forensic content of a different
nature, or scenarios in which investigative relevance depends on contextual
elements not represented in the dataset.

The construction of the corpus also introduces a
\textit{human-in-the-loop} component. This choice increases the semantic quality
of the dataset and reduces the risk of including ambiguous or non-relevant
samples, but it also entails a component of human judgment in the final
selection. This component is neither hidden nor treated as automatic: it is
documented through manifests, inclusion criteria, exclusion criteria, and
balancing. Consequently, reproducibility concerns the frozen experimental
procedure and the final artifacts, whereas the dataset construction phase must
be interpreted as a controlled and auditable process, not as an entirely
automatic procedure.

This experimental scope is also connected to the necessarily selective coverage
of the perturbations and transformations considered. The selected adversarial
attacks and anti-forensic transformations do not cover the entire space of
possible manipulations. The selection was guided by operational
representativeness, computational feasibility, and coherence with the forensic
objective of the thesis, not by exhaustiveness. Model-dependent attacks are
generated against EfficientNet-B0 and must therefore be interpreted in light of
the considered white-box scenario. Anti-forensic transformations, instead, are
model-agnostic and represent operationally plausible image degradations, but
they do not exhaust all possible concealment, manipulation, or falsification
techniques applicable to visual evidence.

An additional element of caution derives from the comparability between proxy
models and selected black-box software systems. Proxy models allow a controlled
evaluation, with access to predictions, confidence, and fully traceable
experimental conditions. Selected software systems, instead, are observable only
through exportable outputs and do not provide access to architecture, training
data, decision thresholds, preprocessing logic, or internal classification
mechanisms. For this reason, the comparison between proxy models and selected
software systems must be understood as an operational triage comparison, not as
an evaluation of homogeneous classifiers on identical labels.

The recoding of black-box outputs into the \texttt{weapon} and
\texttt{non_weapon} classes also introduces a component of interpretive
uncertainty. The adopted mapping is controlled and documented, but it cannot
completely eliminate the limitations arising from the opacity of black-box
systems. This aspect is particularly relevant for
\gls{griffeye}/\gls{t3kcore}, where the observed low \gls{fpr} also reflects
the conservative choice of the primary positive bookmark and cannot be
interpreted as stand-alone evidence of inherently better system performance.

From the perspective of forensic interpretation, it is also necessary to
consider that automatic outputs are not fully separable from the technical
context in which they are produced. The check on sensitive metadata showed that
some selected software outputs may be influenced by information embedded in the
files. The construction of the blind bundle reduces this risk, but it does not
demonstrate the absence of any residual dependency on metadata, internal names,
or non-visual artifacts. This caution is particularly important for selected
black-box software systems, in which it is not possible to directly verify
which signals were used by the system to produce a given classification or
retrieval result.

The \gls{xai} analysis must also be interpreted within a specific scope. The
maps based on \gls{ig} were used as a qualitative study on representative cases
of EfficientNet-B0 and support the diagnostic interpretation of failure cases.
However, they do not constitute a quantitative robustness metric or evidence of
model correctness. Their value lies in making selected local behaviors of the
classifier more directly observable, not in replacing quantitative metrics or
expert validation.

Finally, this work does not include an assessment of the legal admissibility of
the outputs produced by the analyzed \gls{ai}-based systems. The adopted
metrics, such as accuracy, weapon recall, \gls{fnr}, matching rate, and weapon
flag rate, are operational triage metrics and not criteria for evidentiary
validation. The output of a classifier or semantic retrieval system does not, by
itself, constitute an autonomous evidentiary determination: its relevance must
be assessed in the context of the digital chain of custody, the documentation
of the system version and configuration, independent validation, and the
applicable legal framework.

% ─────────────────────────────────────────────────────────────────────────────
\section{Future Work}
\label{sec:conclusions-future-work}
% ─────────────────────────────────────────────────────────────────────────────

Future research directions derive directly from the stated limitations and from
the evidence that emerged throughout the work. The proposed protocol provides a
reproducible basis for evaluating operational robustness, but its value can be
strengthened by extending its experimental scope, improving uncertainty
management, and further consolidating aspects related to verifiability and
auditability.

A natural extension of this work involves applying the protocol to additional
versions, configurations, and operational categories of the selected software
systems analyzed. Since \gls{magnetaxiomtool}/\gls{magnetai},
\gls{excirefoto2025}, \gls{cellebriteinseyets}, and
\gls{griffeye}/\gls{t3kcore} are black-box or near-black-box systems, the
results obtained on a specific version must not be assumed to remain valid for
future releases, different settings, or capabilities not covered in the present
study. A broader validation should therefore repeat the protocol across
software versions, operational configurations, and available semantic
categories, in order to distinguish between stable properties of the system and
behaviors dependent on the specific configuration examined.

The extension of the validation should also concern the dataset and the
experimental tasks. The work could be extended toward multi-class scenarios,
different visual domains, and images acquired under more heterogeneous
operational conditions, including exports from real devices, content originating
from social platforms or messaging applications, and images degraded by forensic
acquisition pipelines. From this perspective, it would be useful to integrate
multi-annotator review procedures or agreement measures on the semantic
selection of samples, in order to quantify the subjective component introduced
by \textit{human-in-the-loop} validation and make the corpus construction
process even more transparent.

The results obtained on \gls{ood} samples also suggest the need to further
investigate uncertainty management. In the experimental protocol, several
systems tend to absorb anomalous content into one of the operational classes,
sometimes with non-negligible confidence levels. Future work could therefore
integrate mechanisms for uncertainty estimation, rejection options, confidence
calibration, and \gls{ood} detection, with the objective of reducing the risk of
forced classifications on inputs that are not relevant to the task. In a
forensic workflow, the ability to recognize when content does not belong to the
expected decision domain is as relevant as the ability to correctly classify
samples belonging to known classes.

With regard to robustness, further developments could consider attack and
defense strategies not included in the present protocol, such as transfer-based
attacks, adversarial patches, physical or hybrid transformations, forms of
adversarial training, and robust preprocessing techniques. These extensions,
however, should remain subordinate to the forensic purpose of the evaluation.
The objective should not be to transform the work into a purely algorithmic
\gls{advml} benchmark, but to understand to what extent specific adversarial or
anti-forensic conditions affect the operational reliability of systems used to
support investigative triage.

A further avenue for development concerns the verifiability of the protocol and
the availability of the artifacts required to reproduce the experiment. Even
when the images cannot be redistributed in full because of licensing constraints
or data sensitivity, scripts, manifests, configurations, hashes, aggregate
metrics, and normalization procedures can be made available to enable
independent verification of the experimental procedure. This requirement is
particularly relevant for selected black-box software systems, for which the
repeatability of the procedure and the clarity of the mapping criteria become
central elements of validation.

Looking ahead, the integration of \gls{xai} methods into forensic software
could also help improve the transparency of automatic outputs. Such integration
should nevertheless be interpreted with caution: visual or textual explanations
can support failure mode diagnosis and critical review by the analyst, but they
must not be treated as stand-alone evidence of classification correctness.
Their value lies in making the behavior of the system more directly observable,
not in replacing quantitative metrics, independent validation, or expert
control.

More generally, this work suggests the need for shared protocols for validating
\gls{ai} in \gls{df}. Such protocols should include metrics oriented toward the
forensic context, software versioning, documentation of configurations, error
analysis, controlled benchmarks, and independent audit procedures. Only through
systematic, traceable, and repeatable validation is it possible to integrate
\gls{ai}-based systems into forensic workflows in a methodologically defensible
manner, while preserving the role of automation as support for triage and not
as a substitute for investigative and expert assessment.

% ─────────────────────────────────────────────────────────────────────────────
\section{Closing Remarks}
\label{sec:conclusions-closing}
% ─────────────────────────────────────────────────────────────────────────────

This thesis has shown that the operational robustness of \gls{ai}-based systems
in \gls{df} cannot be reduced to an accuracy measure on clean data. In the
presence of \gls{ood} inputs, anti-forensic transformations, and adversarial
perturbations, systems with high nominal performance may produce relevant
errors, including high-confidence errors, with direct implications for
investigative triage.

The methodological contribution of this work lies in making this problem
experimentally verifiable and traceable through the \gls{fairlab} protocol and
the forensic evaluation bundle. The transition from an abstract evaluation of
robustness to an evaluation oriented toward the operational conditions of
forensic triage represents the main shift in perspective proposed by the
thesis. From this perspective, robustness is not interpreted as a general and
static property of the model, but as an observable property of the system within
the operational context in which it is used.

Within the considered experimental scope, the results show that none of the
evaluated systems can be considered operationally reliable in an unconditional
sense. Each system exhibits specific failure modes, depending on the input
condition, configuration, task, and type of observable output. This evidence
does not invalidate the use of \gls{ai} in forensic workflows, but defines its
methodological prerequisites: task-specific validation, documentation of the
operational configuration, explicit management of anomalous inputs, artifact
traceability, and human control over the result.

The \gls{fairlab} protocol provides an experimental structure for making these
prerequisites verifiable, auditable, and reproducible. In this sense, the work
does not propose automation as a substitute for the analyst, but as a supporting
component that must be validated, contextualized, and maintained within a
controlled forensic workflow. Operational robustness therefore becomes a
condition to be demonstrated experimentally, not a property to be presumed on
the basis of nominal performance alone.














